\documentclass[12pt]{article}

\usepackage{amsmath,amssymb}
\usepackage{geometry}
\usepackage{setspace}

\geometry{margin=1in}
\setstretch{1.15}

\begin{document}

\begin{center}
\textbf{CSE400 – Fundamentals of Probability in Computing}

\vspace{0.5em}

\textbf{Lecture 5: Bayes’ Theorem, Random Variables, and Probability Mass Function}

\vspace{0.5em}

\textbf{Instructor:} Dhaval Patel, PhD \\
\textbf{Date:} January 20, 2026
\end{center}

\hrule
\vspace{1em}

\section{Bayes’ Theorem}

\subsection{Weighted Average of Conditional Probabilities}

\textbf{Definition and Setup:}  
Let $A$ and $B$ be events. We may express $A$ as
\[
A = AB \cup AB^c
\]
for, in order for an outcome to be in $A$, it must either be in both $A$ and $B$, or be in $A$ but not in $B$.

\textbf{Derivation:}  
Since $AB$ and $AB^c$ are mutually exclusive, by Axiom 3,
\[
\Pr(A) = \Pr(AB) + \Pr(AB^c)
\]
Using the definition of conditional probability,
\[
\Pr(A) = \Pr(A \mid B)\Pr(B) + \Pr(A \mid B^c)\Pr(B^c)
\]
\[
= \Pr(A \mid B)\Pr(B) + \Pr(A \mid B^c)[1 - \Pr(B)]
\]

\textbf{Result:}  
The probability of event $A$ is a weighted average of the conditional probabilities, with weights given by the probability of the event on which it is conditioned.

\subsection{Learning by Example}

\subsubsection{Example 3.1 (Part 1/2)}

\textbf{Problem Statement:}  
An insurance company believes that people can be divided into two classes: those who are accident prone and those who are not.

\begin{itemize}
    \item Probability an accident-prone person has an accident in one year: $( 0.4 )$
    \item Probability a non–accident-prone person has an accident in one year: $( 0.2 )$
    \item Proportion of accident-prone people in the population: $( 0.3 )$
\end{itemize}

What is the probability that a new policyholder will have an accident within a year of purchasing a policy?

\textbf{Solution:}  
Let
\begin{itemize}
    \item $( A_1 )$: event that the policyholder has an accident within a year
    \item $( A )$: event that the policyholder is accident prone
\end{itemize}

Then,
\[
\Pr(A_1) = \Pr(A_1 \mid A)\Pr(A) + \Pr(A_1 \mid A^c)\Pr(A^c)
\]
\[
= (0.4)(0.3) + (0.2)(0.7)
\]
\[
= 0.12 + 0.14 = 0.26
\]

\subsubsection{Example 3.1 (Part 2/2)}

\textbf{Problem Statement:}  
Suppose that a new policyholder has an accident within a year of purchasing a policy. What is the probability that he or she is accident prone?

\textbf{Solution:}  
The desired probability is
\[
\Pr(A \mid A_1) = \frac{\Pr(A \cap A_1)}{\Pr(A_1)}
\]
\[
= \frac{\Pr(A)\Pr(A_1 \mid A)}{\Pr(A_1)}
\]
\[
= \frac{(0.3)(0.4)}{0.26}
\]
\[
= \frac{6}{13}
\]

\subsection{Formal Introduction: Law of Total Probability and Bayes’ Theorem}

\subsubsection{Law of Total Probability (Formula 3.4)}

Let $B_1, B_2, \dots, B_n$ be mutually exclusive and exhaustive events. Then,
\[
\Pr(A) = \sum_{i=1}^{n} \Pr(AB_i)
\]

Using
\[
\Pr(AB_i) = \Pr(B_i \mid A)\Pr(A)
\]

\subsubsection{Bayes’ Formula (Proposition 3.1)}

From the Law of Total Probability,
\[
\Pr(B_i \mid A) = \frac{\Pr(A \mid B_i)\Pr(B_i)}{\sum_{j=1}^{n} \Pr(A \mid B_j)\Pr(B_j)}
\]

\textbf{Terminology:}
\begin{itemize}
    \item $( \Pr(B_i) )$: \textit{a priori probability}
    \item $( \Pr(B_i \mid A) )$: \textit{posterior probability}
\end{itemize}

\subsection{Learning by Example}

\subsubsection{Example 3.2}

\textbf{Problem Statement:}  
There are three cards:
\begin{itemize}
    \item One card is red on both sides (RR)
    \item One card is black on both sides (BB)
    \item One card is red on one side and black on the other (RB)
\end{itemize}

A card is randomly selected and placed on the ground. If the upper side is red, what is the probability that the other side is black?

\textbf{Solution:}  
Let
\begin{itemize}
    \item $( RR, BB, RB )$: events corresponding to the selected card
    \item $( R )$: event that the upturned side is red
\end{itemize}

Then,
\[
\Pr(RB \mid R) = \frac{\Pr(R \mid RB)\Pr(RB)}{\Pr(R \mid RR)\Pr(RR) + \Pr(R \mid RB)\Pr(RB) + \Pr(R \mid BB)\Pr(BB)}
\]

Substituting values,
\[
= \frac{(1/2)(1/3)}{(1)(1/3) + (1/2)(1/3) + (0)(1/3)}
\]
\[
= \frac{1/6}{1/2} = \frac{1}{3}
\]

\section{Random Variables}

\subsection{Motivation and Concept}

When an experiment is performed, we are frequently interested mainly in some function of the outcome rather than the outcome itself.

Examples described in lecture:
\begin{itemize}
    \item Dice tossing: interest in the sum of two dice
    \item Coin flipping: interest in the total number of heads
\end{itemize}

\textbf{Definition:}  
A \textit{random variable} is a real-valued function defined on the sample space.

\begin{itemize}
    \item Values are determined by outcomes of an experiment
    \item Probabilities are assigned to possible values of random variables
\end{itemize}

The distribution of a random variable can be visualized as a bar diagram where:
\begin{itemize}
    \item x-axis represents values of the random variable
    \item Height of the bar at value $( a )$ is $( \Pr[X = a] )$
\end{itemize}

\subsection{Examples}

\subsubsection{Example 1: Tossing 3 Fair Coins}

\textbf{Problem Statement:}  
Suppose an experiment consists of tossing 3 fair coins. Let $( Y )$ denote the number of heads that appear.

\textbf{Possible Values:}
\[
Y \in \{0, 1, 2, 3\}
\]

\textbf{Probabilities:}
\[
\Pr(Y = 0) = \Pr\{(t,t,t)\} = \frac{1}{8}
\]
\[
\Pr(Y = 1) = \Pr\{(t,t,h),(t,h,t),(h,t,t)\} = \frac{3}{8}
\]
\[
\Pr(Y = 2) = \Pr\{(t,h,h),(h,t,h),(h,h,t)\} = \frac{3}{8}
\]
\[
\Pr(Y = 3) = \Pr\{(h,h,h)\} = \frac{1}{8}
\]

Since $( Y )$ must take one of these values,
\[
\sum_{y=0}^{3} \Pr(Y = y) = 1
\]

\section{Probability Mass Function (PMF)}

\subsection{Concept}

A random variable that can take on at most a countable number of possible values is said to be \textit{discrete}.

Let $( X )$ be a discrete random variable with range
\[
R_X = \{x_1, x_2, x_3, \dots\}
\]
(finite or countably infinite).

\textbf{Definition:}  
The function
\[
p_X(x_k) = \Pr(X = x_k)
\]
is called the \textit{Probability Mass Function (PMF)} of $( X )$.

Since $( X )$ must take on one of the values $( x_k )$,
\[
\sum_k p_X(x_k) = 1
\]

\subsection{Example: Two Independent Tosses of a Fair Coin}

(As presented in lecture slides.)

\subsection{Example}

\textbf{Problem Statement:}  
The probability mass function of a random variable $( X )$ is given by
\[
p(i) = c \lambda^i, \quad i = 0,1,2,\dots
\]
where $( \lambda > 0 )$. Find:
\begin{itemize}
    \item $( \Pr(X = 0) )$
    \item $( \Pr(X > 2) )$
\end{itemize}

\textbf{Solution:}  
Since $( p(i) )$ is a PMF,
\[
\sum_{i=0}^{\infty} c \lambda^i = 1
\]

Using the geometric series,
\[
c \sum_{i=0}^{\infty} \lambda^i = c \cdot \frac{1}{1 - \lambda} = 1
\]
\[
\Rightarrow c = 1 - \lambda
\]

Then,
\[
\Pr(X = 0) = p(0) = c = 1 - \lambda
\]

Also,
\[
\Pr(X > 2) = 1 - \Pr(X \le 2)
\]
\[
= 1 - [p(0) + p(1) + p(2)]
\]

\hrule
\begin{center}
\textbf{End of Lecture 5 Scribe}
\end{center}

\end{document}
